\section{Architectural and Integration Patterns}
\label{sec:patterns}

The solution combines four architectural patterns with six enterprise integration patterns. Together they enable heterogeneous protocol integration, high-volume asynchronous processing, real-time push notifications, distributed transaction management, and independent scalability.

\subsection{Architectural Patterns}
\label{subsec:arch_patterns}

\subsubsection{Microservices Architecture}

Each business capability is deployed as a separately runnable and independently scalable unit. The thirteen containers are grouped by role: client-facing services (Auth, Order, Items, Delivery, Notify), backend processing services (Payment, Route), legacy integration adapters (CMS, ROS, WMS), and infrastructure (RabbitMQ, PostgreSQL). Because each container has a single responsibility, the CMS Adapter can be redeployed when the SOAP contract changes without touching the Order or Payment services. In a Kubernetes deployment, only the services under highest load need to be scaled out, leaving all others at a single replica.

\begin{figure}[H]
    \centering
    \includegraphics[width=0.9\textwidth]{figures/microservices_pattern.pdf}
    \caption{Microservices Architecture Pattern}
    \label{fig:microservices}
\end{figure}

\subsubsection{Event-Driven Architecture}

Services communicate by publishing and consuming events rather than making direct synchronous calls to each other. Every state change in the order lifecycle is expressed as a routing key published to the \texttt{swift\_logistics} topic exchange: \texttt{order.created}, \texttt{order.confirmed}, \texttt{order.wms.reserved}, \texttt{order.payment.completed}, and others. Producers have no knowledge of their consumers. Adding a new service that must react to an existing event requires only a new queue binding, with no change to any existing service. Routing key definitions are declared in each service's \texttt{config.js} under \texttt{publishedRoutingKeys} and \texttt{subscribedRoutingKeys}.

\begin{figure}[H]
    \centering
    \includegraphics[width=0.9\textwidth]{figures/event_driven_pattern.pdf}
    \caption{Event-Driven Architecture Pattern}
    \label{fig:event_driven}
\end{figure}

\subsubsection{API Gateway Pattern}

Implemented in \texttt{source-code/services/api-gateway/index.js}. The gateway handles three cross-cutting concerns before forwarding any request: authentication (\texttt{validateToken()} verifies the HS256 JWT and rejects with HTTP~401), identity propagation (verified \texttt{username} and \texttt{role} are injected as \texttt{x-username} and \texttt{x-role} headers), and routing (\texttt{resolveTarget()} maps URL prefix to internal container address). WebSocket upgrade requests are proxied via the \texttt{upgrade} event handler. Internal services are entirely hidden from external network interfaces; the gateway is the sole entry point.

\subsubsection{Service Mesh Pattern}

A service mesh is not implemented in the prototype. Docker's internal bridge network and direct AMQP connections provide network-level isolation at prototype scale. In production, a mesh such as Istio or Linkerd would add mutual TLS between containers, distributed tracing, and traffic shaping without requiring code changes to individual services.

\subsection{Integration Patterns}
\label{subsec:integration_patterns}

\subsubsection{Message Channel}

RabbitMQ queues are the message channels. All queues are declared with \texttt{durable: true}, so messages survive broker restarts. Two channel types are in use. Point-to-point: each application service has a stable named queue (e.g.\ \texttt{order-service-queue.*}) consumed only by its own instances, ensuring each message is processed exactly once. Fanout: the Notify Service declares a unique queue per instance (\texttt{notifier.\$\{HOSTNAME\}}) bound to \texttt{notify.*}, so every running instance receives every notification event independently. Queue declaration and binding is handled by the shared \texttt{utility/pubsub.js} utility.

\begin{figure}[H]
    \centering
    \includegraphics[width=0.9\textwidth]{figures/message_channel_pattern.pdf}
    \caption{Message Channel Pattern}
    \label{fig:message_channel}
\end{figure}

\subsubsection{Message Translator}

Each adapter acts as a Message Translator between the canonical JSON used by the application layer and the native protocol of its legacy system. The CMS Adapter receives \texttt{\{correlationId, clientId, orderData\}} and translates it into a SOAP/XML \texttt{CreateOrder} envelope using the \texttt{soap} library. The ROS Adapter receives \texttt{\{correlationId, locations[]\}} and translates it to an HTTP POST to \texttt{/api/ros/optimize-route} via \texttt{axios}. The WMS Adapter receives \texttt{\{correlationId, orderId, itemList\}} and translates it to \texttt{POST /api/wms/reserve}. All three adapters present an identical AMQP interface to the rest of the system regardless of the underlying protocol. Source: \texttt{cms-adapter/index.js}, \texttt{ros-adapter/index.js}, \texttt{wms-adapter/index.js}.

\begin{figure}[H]
    \centering
    \includegraphics[width=0.9\textwidth]{figures/message_translator_pattern.pdf}
    \caption{Message Translator Pattern}
    \label{fig:translator}
\end{figure}

\subsubsection{Content-Based Router}

RabbitMQ's topic exchange acts as the content-based router: the routing key attached to each published message determines which queues receive it. When the Order Service publishes \texttt{order.confirmed}, only the WMS Adapter's queue (bound to \texttt{order.confirmed}) receives it; CMS and ROS adapters are not notified. The \texttt{notify.*} wildcard pattern routes all \texttt{notify.order.*} and \texttt{notify.client.*} events to all Notify Service instance queues. No application code is required to implement this routing; it is entirely driven by exchange type and binding keys.

\subsubsection{Request-Reply (Async RPC) Pattern}

Used wherever a service needs a response from a legacy system without blocking the HTTP thread. The Items Service publishes \texttt{wms.items.request} with a \texttt{correlationId}, stores the ID in a pending-requests map, and returns HTTP~200 immediately. When the WMS Adapter later publishes \texttt{wms.items.response}, the Notify Service delivers the result to the client via WebSocket. The Order Service uses the same pattern for order creation: it publishes \texttt{cms.order.create} with the \texttt{x-request-id} idempotency header as the \texttt{correlationId} and awaits the CMS reply asynchronously, ensuring exactly-once processing for duplicate requests.

\subsubsection{Publish-Subscribe Pattern}

Any service can broadcast an event to multiple independent consumers by publishing a routing key to the \texttt{swift\_logistics} exchange, which routes it to all matching queues. A key example is \texttt{order.confirmed}: it is consumed by the WMS Adapter to reserve stock. Future consumers can subscribe to the same key without any change to the publisher. The Notify Service's \texttt{notify.*} subscription is the primary real-time delivery mechanism: any service publishes \texttt{notify.<userId>.<event>} and the Notify Service pushes the payload over WebSocket to the connected user.

\begin{figure}[H]
    \centering
    \includegraphics[width=0.9\textwidth]{figures/pubsub_pattern.pdf}
    \caption{Publish-Subscribe Pattern}
    \label{fig:pubsub}
\end{figure}

\subsubsection{Competing Consumers}

Multiple consumer instances can share a single durable queue; RabbitMQ delivers each message to exactly one instance. The prototype runs one instance per service, but the design natively supports scale-out. Adding a second \texttt{order-service} container means both instances consume from the same queue set and each message is handled exactly once. The Notify Service is the exception: it uses unique per-instance queues so every instance independently receives every notification event and serves its own connected users.

\subsubsection{Dead Letter Channel}

The Dead Letter Channel pattern is architecturally designed but not implemented in the prototype. In production, each adapter queue would be configured with an \texttt{x-dead-letter-exchange} and \texttt{x-message-ttl}. If an adapter cannot process a message and nacks it without requeueing, RabbitMQ would route it to a dead letter exchange where a monitoring consumer could log the failure and schedule a retry. Currently, adapters log exceptions but ACK the message regardless, meaning persistently failing messages are silently dropped. This is a known limitation noted in Section~\ref{subsec:limitations}.

\subsubsection{Saga Pattern (Choreography)}

The order creation saga is a sequence of local transactions linked by events. No central orchestrator exists; each service reacts to an incoming event and emits the next event or a compensating event. The Order Service tracks in-flight state in a \texttt{pendingPayments} Map keyed by \texttt{correlationId}.

The nine-step happy path is: (1)~\texttt{POST /order}~\textrightarrow~Order Service publishes \texttt{cms.order.create}; (2)~CMS Adapter creates the order via SOAP~\textrightarrow~publishes \texttt{order.created}; (3)~Payment Service creates a payment intent~\textrightarrow~notifies client; (4)~Client confirms via \texttt{PATCH /order}~\textrightarrow~Order Service publishes \texttt{cms.order.update\_status}; (5)~CMS Adapter updates status~\textrightarrow~publishes \texttt{order.status\_updated}; (6)~Order Service publishes \texttt{order.confirmed}; (7)~WMS Adapter reserves stock~\textrightarrow~publishes \texttt{order.wms.reserved}; (8)~Order Service publishes \texttt{payment.charge}; (9a)~Payment succeeds~\textrightarrow~publishes \texttt{order.payment.completed}.

On payment failure (9b), the Payment Service publishes \texttt{order.payment.failed}. The Order Service publishes \texttt{wms.order.release} as the compensating action; the WMS Adapter releases the reservation. Source: \texttt{order-service/index.js}, \texttt{payment-service/index.js}, \texttt{wms-adapter/index.js}.

\begin{figure}[H]
    \centering
    \includegraphics[width=\textwidth]{figures/saga_pattern.pdf}
    \caption{Saga Pattern for Distributed Transactions}
    \label{fig:saga}
\end{figure}

\subsection{Pattern Interaction and Composition}
\label{subsec:pattern_composition}

\subsubsection{Order Submission Flow}

The order submission flow combines five patterns. The API Gateway Pattern handles authentication and routing of \texttt{POST /order}. The Async RPC pattern (\texttt{correlationId}) links the Order Service to the CMS Adapter, which uses the Message Translator pattern to convert the request into a SOAP envelope. The Publish-Subscribe pattern broadcasts \texttt{order.confirmed} to the WMS Adapter. The Saga pattern manages the distributed transaction across CMS, WMS, and Payment Services. Finally, the Notify Service delivers the payment outcome to the client via WebSocket. The full workflow is documented in \texttt{workflow-diagrams/order-workflow.mmd}.

\begin{figure}[H]
    \centering
    \includegraphics[width=\textwidth]{figures/order_flow_sequence.pdf}
    \caption{Order Submission Sequence Diagram}
    \label{fig:order_sequence}
\end{figure}

\subsubsection{Real-time Update Flow}

When a driver marks a delivery as completed (\texttt{PATCH /deliveries}), the Delivery Service publishes \texttt{wms.delivery.update\_status} (Publish-Subscribe). The WMS Adapter translates the event to an HTTP PATCH on the mock WMS (Message Translator) and replies on a reply key (Request-Reply). The Delivery Service receives confirmation and publishes \texttt{notify.<clientId>.delivery\_update} (Publish-Subscribe). The Notify Service pushes the frame over WebSocket to the client if online, or stores it in \texttt{pending\_notifications} if offline. The full flow is captured in \texttt{workflow-diagrams/delivery-complete-workflow.mmd}.

\begin{figure}[H]
    \centering
    \includegraphics[width=\textwidth]{figures/realtime_update_sequence.pdf}
    \caption{Real-time Update Sequence Diagram}
    \label{fig:realtime_sequence}
\end{figure}

\subsection{Rationale for Pattern Selection}
\label{subsec:pattern_rationale}

Table~\ref{tab:pattern_challenge} maps each architectural challenge to the patterns that address it.

\begin{table}[H]
    \centering
    \caption{Pattern-to-Challenge Mapping}
    \label{tab:pattern_challenge}
    \renewcommand{\arraystretch}{1.2}
    \begin{tabular}{|p{4.5cm}|p{4.5cm}|p{5cm}|}
        \hline
        \textbf{Challenge} & \textbf{Pattern(s)} & \textbf{Justification} \\
        \hline
        Heterogeneous Integration & Message Translator & Encapsulates protocol conversion in one component per system. \\
        \hline
        Real-time Tracking & Pub-Sub + EDA & Notify Service subscribes to \texttt{notify.*} events and pushes over WebSocket without polling. \\
        \hline
        High-Volume Processing & Message Channel + Competing Consumers & Durable queues buffer orders during traffic spikes; multiple consumer instances drain the queue in parallel. \\
        \hline
        Transaction Management & Saga (choreography) + Async RPC & Each step is an independent local transaction; \texttt{correlationId} links steps; compensating events provide rollback. \\
        \hline
        Scalability and Resilience & Microservices + Competing Consumers & Each service scales independently; durable queues prevent data loss when a consumer restarts. \\
        \hline
        Security & API Gateway Pattern & Single JWT enforcement point; no internal service is directly reachable from outside the Docker network. \\
        \hline
    \end{tabular}
\end{table}
